\chapter*{Global Exercise Bank --- Linked Comprehension Passages}
\addcontentsline{toc}{chapter}{Global Exercise Bank: Linked Comprehension Passages}
\label{chap:comprehension}

\begin{tcolorbox}[enhanced,colback=subtleblue,colframe=primarynavy,arc=2mm,boxrule=1pt,
    title=\textbf{\large Linked Comprehension / Case-Study Questions (All Chapters)}]
Paragraph-based and scenario-linked problem sets with multi-step analytical reasoning.
Sources: \textbf{[NA]} Level-3 Comprehensions, \textbf{[NRJ]} Ex-II Comprehension, \textbf{[PRSN]} Advanced Passages.
\end{tcolorbox}

\section*{Passage 1: Debye--H\"{u}ckel--Onsager Theory and High-Field Effects}

The Debye--H\"{u}ckel--Onsager (DHO) equation for a $1:1$ univalent electrolyte in dilute aqueous solution at $25^\circ\text{C}$ is given by:
\begin{equation*}
    \Lambda_m = \Lambda_m^\circ - (A + B\Lambda_m^\circ)\sqrt{c}
\end{equation*}
where $A$ accounts for the electrophoretic drag of the solvent countercurrent and $B$ accounts for the relaxation (asymmetry) field retardance caused by the finite relaxation time $\tau$ of the ionic atmosphere. Under ordinary conditions, the external electric field is small ($E \sim 1\text{--}10\,\text{V\,cm}^{-1}$), so the velocity of ion migration is small compared to the thermal motion, allowing the ionic atmosphere to be treated as perturbed linearly.

However, when an extremely high potential gradient is applied ($E > 10^5\,\text{V\,cm}^{-1}$), an ion moves so rapidly that it travels a distance greater than the thickness of the ionic cloud ($1/\kappa$) within the relaxation time $\tau$. Consequently, the ion escapes its retarding counter-ion cloud, and both asymmetry and electrophoretic effects diminish to zero (the \textbf{Wien effect}). Furthermore, at very high alternating frequencies ($\nu > 10^6\,\text{Hz}$), the direction of the field alternates faster than the relaxation time $\tau$, so the ionic atmosphere remains essentially spherically symmetric, yielding an increased conductivity (the \textbf{Debye--Falkenhagen effect}).

\begin{problembox}
\textbf{CP.1.1} \hfill \textbf{[NA Level-3 / ATK]}

Which of the following statements correctly explains the Wien effect?
\begin{enumerate}[label=(\alph*)]
    \item High electric field increases the thermal ionization energy of solvent molecules.
    \item Under high electric field gradient, an ion traverses a distance greater than the radius of its ionic atmosphere within the relaxation time, stripping off the cloud.
    \item High voltage destroys the hydration shell of the ions permanently.
    \item Viscosity of water drops to zero at high electric potential gradients.
\end{enumerate}
\end{problembox}
\begin{solution}
\textbf{(b)} In very high electric fields ($> 10^4\text{--}10^5\,\text{V\,cm}^{-1}$), the ion velocity $v = uE$ becomes so great that the ion migrates out of its ionic atmosphere before the atmosphere can reform behind it. The asymmetric retarding force and electrophoretic drag vanish, causing $\Lambda_m$ to approach its limiting value $\Lambda_m^\circ$ even at finite concentration.
\end{solution}

\begin{problembox}
\textbf{CP.1.2} \hfill \textbf{[ATK / PRSN]}

The Debye--Falkenhagen effect predicts that when the frequency of the alternating current is increased beyond $\sim 10\,\text{MHz}$, the molar conductance of a strong electrolyte solution:
\begin{enumerate}[label=(\alph*)]
    \item Decreases to zero because ions cannot follow the field oscillations.
    \item Remains unchanged because relaxation time is independent of frequency.
    \item Increases because the ionic atmosphere has no time to form an asymmetric retardance behind the oscillating ion.
    \item First decreases and then approaches infinity.
\end{enumerate}
\end{problembox}
\begin{solution}
\textbf{(c)} At radio frequencies ($\sim 10^7\,\text{Hz}$), the period of oscillation is shorter than the relaxation time of the ionic atmosphere ($\tau \approx 10^{-8}\text{--}10^{-10}\,\text{s}$). As a result, the ionic atmosphere never has time to become asymmetric with respect to the central ion; it remains spherically symmetric, and the relaxation retarding force vanishes, leading to an increase in conductance.
\end{solution}

\begin{problembox}
\textbf{CP.1.3} \hfill \textbf{[NRJ Ex-II / NA]}

For a $0.0004\,\text{M}$ solution of \ce{KCl} at $25^\circ\text{C}$, $\Lambda_m = 148.2\,\text{S\,cm}^2\,\text{mol}^{-1}$ and $\Lambda_m^\circ = 149.8\,\text{S\,cm}^2\,\text{mol}^{-1}$. The experimental Onsager slope $b_{\text{obs}}$ in $\text{S\,cm}^2\,\text{mol}^{-1}\,(\text{mol\,L}^{-1})^{-1/2}$ is:
\begin{multicols}{4}
\begin{enumerate}[label=(\alph*)]
    \item $80$
    \item $40$
    \item $160$
    \item $0.80$
\end{enumerate}
\end{multicols}
\end{problembox}
\begin{solution}
\textbf{(a)}
Using the Onsager equation $\Lambda_m = \Lambda_m^\circ - b\sqrt{c}$:
\begin{align*}
    148.2 &= 149.8 - b\sqrt{0.0004} = 149.8 - b(0.02) \\
    0.02 b &= 149.8 - 148.2 = 1.6 \\
    b &= \frac{1.6}{0.02} = 80\,\text{S\,cm}^2\,\text{mol}^{-1/2}\,\text{L}^{1/2}
\end{align*}
Hence option (a) is correct.
\end{solution}

\vspace{0.5cm}
\section*{Passage 2: Concentration Cells and Liquid Junction Potential}

A concentration cell consists of two half-cells with identical electrodes and electrolytes but differing electrolyte activities or electrode pressures. Concentration cells can be classified into:
\begin{itemize}
    \item \textbf{Electrode concentration cells:} Differing gas pressures or amalgam activities with uniform solution.
    \item \textbf{Electrolyte concentration cells:} Identical electrodes dipping into solutions of the same electrolyte at different mean activities $a_1$ and $a_2$ ($a_2 > a_1$).
\end{itemize}

When the two solutions of different concentrations are in direct liquid-liquid contact without a salt bridge, a liquid junction potential ($E_{lj}$) develops at the interface due to the unequal diffusion rates of cations and anions. For a $1:1$ uni-univalent electrolyte (\ce{M+A-}), the EMF with transference ($E_{wt}$) and without transference ($E_{wot}$) are given by:
\begin{equation*}
    E_{wt} = 2 t_- \frac{RT}{F} \ln \frac{a_2}{a_1}, \qquad E_{wot} = \frac{RT}{F} \ln \frac{a_2}{a_1}
\end{equation*}
where $t_-$ is the transport number of the anion. The liquid junction potential is defined as $E_{lj} = E_{wt} - E_{wot} = (2t_- - 1)\frac{RT}{F}\ln \frac{a_2}{a_1}$. A salt bridge containing \ce{KCl} or \ce{NH4NO3} in agar-agar gel is employed because $t_+ \approx t_- \approx 0.5$ for these electrolytes, practically nullifying $E_{lj}$.

\begin{problembox}
\textbf{CP.2.1} \hfill \textbf{[ATK Topic 6C / NA]}

Consider the cell: $\ce{Pt} \mid \ce{H2}(1\,\text{atm}) \mid \ce{HCl}(a_1) \parallel \ce{HCl}(a_2) \mid \ce{H2}(1\,\text{atm}) \mid \ce{Pt}$, where a salt bridge eliminates the liquid junction. If $a_2 = 10\,a_1$ at $298\,\text{K}$, the cell EMF is:
\begin{multicols}{4}
\begin{enumerate}[label=(\alph*)]
    \item $0.0591\,\text{V}$
    \item $0.1182\,\text{V}$
    \item $0.0296\,\text{V}$
    \item $0\,\text{V}$
\end{enumerate}
\end{multicols}
\end{problembox}
\begin{solution}
\textbf{(a)}
The cell without transference:
\begin{align*}
    \text{Anode: } & \tfrac{1}{2}\ce{H2} \to \ce{H+}(a_1) + e^- \\
    \text{Cathode: } & \ce{H+}(a_2) + e^- \to \tfrac{1}{2}\ce{H2} \\
    \text{Overall: } & \ce{H+}(a_2) \to \ce{H+}(a_1)
\end{align*}
\begin{equation*}
    E_{\text{cell}} = 0 - \frac{RT}{F}\ln \frac{a_1}{a_2} = \frac{2.303 RT}{F} \log_{10} \frac{a_2}{a_1} = 0.0591 \times \log_{10}(10) = 0.0591\,\text{V}
\end{equation*}
\end{solution}

\begin{problembox}
\textbf{CP.2.2} \hfill \textbf{[PRSN / NRJ]}

If the cell with direct boundary $\ce{Pt, H2}(1\,\text{atm}) \mid \ce{HCl}(a_1) : \ce{HCl}(a_2) \mid \ce{H2}(1\,\text{atm}), \ce{Pt}$ is constructed without a salt bridge, and the transport number of $\ce{H+}$ is $t_+ = 0.82$, then the liquid junction potential $E_{lj}$ for $a_2/a_1 = 10$ at $25^\circ\text{C}$ is:
\begin{multicols}{4}
\begin{enumerate}[label=(\alph*)]
    \item $+0.0378\,\text{V}$
    \item $-0.0378\,\text{V}$
    \item $+0.0591\,\text{V}$
    \item $-0.0189\,\text{V}$
\end{enumerate}
\end{multicols}
\end{problembox}
\begin{solution}
\textbf{(b)}
Given $t_+ = 0.82$, thus $t_- = 1 - t_+ = 0.18$.
\begin{equation*}
    E_{lj} = (2t_- - 1)\frac{2.303 RT}{F}\log_{10} \frac{a_2}{a_1} = [2(0.18) - 1](0.0591)\log_{10}(10) = (0.36 - 1)(0.0591) = -0.64 \times 0.0591 = -0.0378\,\text{V}
\end{equation*}
The negative sign indicates that the more dilute solution develops a positive potential relative to the concentrated solution due to faster diffusion of $\ce{H+}$ ions across the boundary.
\end{solution}

\begin{problembox}
\textbf{CP.2.3} \hfill \textbf{[NA Level-3 / CNG-TH]}

Why is \ce{KNO3} or \ce{NH4NO3} chosen over \ce{NaCl} for salt bridges in general electrochemistry?
\begin{enumerate}[label=(\alph*)]
    \item \ce{Na+} reacts with most metal cations to precipitate hydroxides.
    \item \ce{K+} and \ce{NO3-} have nearly equal ionic mobilities, so $t_+ \approx t_- \approx 0.5$, minimizing $E_{lj}$ to under $1\text{--}2\,\text{mV}$.
    \item \ce{NaCl} is insoluble in agar-agar gel.
    \item Nitrates undergo redox oxidation at the cathode during discharge.
\end{enumerate}
\end{problembox}
\begin{solution}
\textbf{(b)} For a salt bridge electrolyte to effectively eliminate the liquid junction potential, the transport numbers of the cation and anion must be almost identical ($t_+ \approx t_- \approx 0.5$). In \ce{KNO3}, $u(\ce{K+}) \approx 7.6 \times 10^{-4}$ and $u(\ce{NO3-}) \approx 7.4 \times 10^{-4}\,\text{cm}^2\,\text{V}^{-1}\,\text{s}^{-1}$, yielding $E_{lj} \approx 0$. In contrast, for \ce{NaCl}, $u(\ce{Cl-}) \gg u(\ce{Na+})$ ($t_+ \approx 0.39, t_- \approx 0.61$), generating a substantial non-zero junction potential.
\end{solution}

\vspace{0.5cm}
\section*{Passage 3: Thermodynamics of the Lead-Acid Accumulator}

The lead-acid storage battery is the ubiquitous secondary electrochemical energy source in automobiles. Its cell representation in the fully charged state is:
\begin{equation*}
    \ce{Pb(s)} \mid \ce{PbSO4(s)} \mid \ce{H2SO4(aq, } 38\%\, w/w, \rho \approx 1.28\,\text{g\,cm}^{-3}\ce{)} \mid \ce{PbSO4(s)} \mid \ce{PbO2(s)} \mid \ce{Pb(s)}
\end{equation*}
During discharge, the electrode half-reactions and overall reversible cell process are:
\begin{align*}
    \text{Anode: } & \ce{Pb(s) + HSO4^-(aq) -> PbSO4(s) + H+(aq) + 2e^-} \\
    \text{Cathode: } & \ce{PbO2(s) + 3H+(aq) + HSO4^-(aq) + 2e^- -> PbSO4(s) + 2H2O(l)} \\
    \text{Net Cell Reaction: } & \ce{Pb(s) + PbO2(s) + 2H2SO4(aq) -> 2PbSO4(s) + 2H2O(l)}
\end{align*}
For this reaction, $n = 2$, $E^\circ_{\text{cell}} \approx 2.05\,\text{V}$ at $298\,\text{K}$, and $\left(\frac{\partial E}{\partial T}\right)_P \approx +1.5 \times 10^{-4}\,\text{V\,K}^{-1}$. During continuous discharge, \ce{H2SO4} is consumed and \ce{H2O} is produced, causing the electrolyte density to fall towards $\rho \approx 1.15\,\text{g\,cm}^{-3}$, providing an accurate measure of the battery's state of charge via a hydrometer.

\begin{problembox}
\textbf{CP.3.1} \hfill \textbf{[GRB / PRSN]}

What is the entropy change $\Delta S$ for the discharge of the lead-acid cell per mole of \ce{Pb} consumed at $298\,\text{K}$?
\begin{multicols}{4}
\begin{enumerate}[label=(\alph*)]
    \item $+28.95\,\text{J\,K}^{-1}$
    \item $-28.95\,\text{J\,K}^{-1}$
    \item $+14.48\,\text{J\,K}^{-1}$
    \item $+57.90\,\text{J\,K}^{-1}$
\end{enumerate}
\end{multicols}
\end{problembox}
\begin{solution}
\textbf{(a)}
From the thermodynamic relationship $\Delta S = nF \left(\frac{\partial E}{\partial T}\right)_P$:
\begin{equation*}
    \Delta S = 2 \times 96500\,\text{C\,mol}^{-1} \times (+1.5 \times 10^{-4}\,\text{V\,K}^{-1}) = +28.95\,\text{J\,K}^{-1}\,\text{mol}^{-1}
\end{equation*}
Since $\Delta S > 0$, the cell absorbs heat from the surroundings ($q_{\text{rev}} = T\Delta S > 0$) during reversible discharge at constant temperature.
\end{solution}

\begin{problembox}
\textbf{CP.3.2} \hfill \textbf{[NA Level-3 / NRJ]}

A lead-acid battery delivers a steady current of $5.0\,\text{A}$ for $10.0\,\text{hours}$. Assuming $100\%$ current efficiency, the total mass of \ce{H2SO4} consumed from the electrolyte is: ($M_{\text{H2SO4}} = 98.08\,\text{g\,mol}^{-1}$)
\begin{multicols}{4}
\begin{enumerate}[label=(\alph*)]
    \item $183\,\text{g}$
    \item $91.5\,\text{g}$
    \item $366\,\text{g}$
    \item $245\,\text{g}$
\end{enumerate}
\end{multicols}
\end{problembox}
\begin{solution}
\textbf{(a)}
Total charge passed:
\begin{equation*}
    Q = I \times t = 5.0\,\text{A} \times (10 \times 3600\,\text{s}) = 180000\,\text{C}
\end{equation*}
Number of moles of electrons transferred:
\begin{equation*}
    n_e = \frac{Q}{F} = \frac{180000}{96500} \approx 1.8653\,\text{mol of } e^-
\end{equation*}
From the balanced net reaction, $2\,\text{mol of } e^-$ correspond to consumption of $2\,\text{mol of } \ce{H2SO4}$.
Thus, moles of \ce{H2SO4} consumed = moles of $e^- = 1.8653\,\text{mol}$.
Mass of \ce{H2SO4} consumed:
\begin{equation*}
    m = 1.8653\,\text{mol} \times 98.08\,\text{g\,mol}^{-1} \approx 182.95\,\text{g} \approx 183\,\text{g}
\end{equation*}
\end{solution}

\begin{problembox}
\textbf{CP.3.3} \hfill \textbf{[CNG-TH / GRB]}

During the recharging process of a discharged lead storage cell, which of the following processes occurs at the anode (electrode connected to the positive terminal of the external DC source)?
\begin{enumerate}[label=(\alph*)]
    \item \ce{PbSO4(s)} is reduced to spongy \ce{Pb(s)}.
    \item \ce{PbSO4(s)} is oxidized to \ce{PbO2(s)}.
    \item \ce{H2SO4} is consumed, lowering the electrolyte density further.
    \item Hydrogen gas is evolved briskly from the \ce{PbO2} grid.
\end{enumerate}
\end{problembox}
\begin{solution}
\textbf{(b)} During recharging, the cell functions as an electrolytic cell. The positive plate (anode) undergoes oxidation:
\begin{equation*}
    \ce{PbSO4(s) + 2H2O(l) -> PbO2(s) + 3H+(aq) + HSO4^-(aq) + 2e^-}
\end{equation*}
Meanwhile, reduction occurs at the cathode (negative terminal): $\ce{PbSO4(s) + H+(aq) + 2e^- -> Pb(s) + HSO4^-(aq)}$. Sulfuric acid is regenerated, raising the density.
\end{solution}

\vspace{0.5cm}
\section*{Passage 4: Frost and Latimer Diagrams of Chlorine}

The standard reduction potentials (in Volts) for various chlorine species in acidic aqueous solution ($\text{pH} = 0$) are represented by the following Latimer diagram:
\begin{equation*}
    \ce{ClO4^-} \xrightarrow{+1.201\,\text{V}} \ce{ClO3^-} \xrightarrow{+1.181\,\text{V}} \ce{HClO2} \xrightarrow{+1.674\,\text{V}} \ce{HClO} \xrightarrow{+1.630\,\text{V}} \ce{Cl2} \xrightarrow{+1.358\,\text{V}} \ce{Cl^-}
\end{equation*}
For any half-reaction connecting oxidation state $N_1$ to $N_2$ ($\Delta N = N_1 - N_2$), the standard Gibbs energy is $\Delta G^\circ = -\Delta N \cdot F \cdot E^\circ$. On a Frost diagram, the quantity $N \cdot E^\circ$ (which is proportional to $\Delta G^\circ / F$ relative to the element in its standard state $N=0$) is plotted on the vertical axis against the oxidation number $N$ on the horizontal axis. A species is thermodynamically unstable with respect to disproportionation if it lies \textit{above} the chord connecting its two adjacent oxidation states, and stable against disproportionation if it lies \textit{below} that chord.

\begin{problembox}
\textbf{CP.4.1} \hfill \textbf{[ATK Topic 6B / NA]}

What is the standard reduction potential $E^\circ$ for the direct conversion of chlorate to chlorine gas:
\begin{equation*}
    \ce{2ClO3^- + 12H+ + 10e^- -> Cl2 + 6H2O}?
\end{equation*}
\begin{multicols}{4}
\begin{enumerate}[label=(\alph*)]
    \item $+1.47\,\text{V}$
    \item $+1.29\,\text{V}$
    \item $+1.38\,\text{V}$
    \item $+1.61\,\text{V}$
\end{enumerate}
\end{multicols}
\end{problembox}
\begin{solution}
\textbf{(a)}
The steps from $\ce{ClO3^-}$ ($N = +5$) to $\ce{Cl2}$ ($N = 0$) involve $n = 5$ electrons per Cl atom:
\begin{align*}
    \ce{ClO3^- -> HClO2} \quad &(n_1 = 1, E_1^\circ = 1.181\,\text{V}) \implies n_1 E_1^\circ = 1.181\,\text{V} \\
    \ce{HClO2 -> HClO} \quad &(n_2 = 1, E_2^\circ = 1.674\,\text{V}) \implies n_2 E_2^\circ = 1.674\,\text{V} \\
    \ce{HClO -> Cl2} \quad &(n_3 = 1, E_3^\circ = 1.630\,\text{V}) \implies n_3 E_3^\circ = 1.630\,\text{V} \\
    \ce{Cl2 -> Cl^-} \quad &\text{(not involved here)}
\end{align*}
Summing $\sum n_i$: $n = 1 + 1 + 1 = 3$ (wait: from $+5$ to $+3$ is $2e^-$, let's verify oxidation states: $\ce{ClO3^-} (+5)$, $\ce{HClO2} (+3)$, $\ce{HClO} (+1)$, $\ce{Cl2} (0)$).
Ah! $\Delta N$ for $\ce{ClO3^- -> HClO2}$ is $5 - 3 = 2e^-$.
$\Delta N$ for $\ce{HClO2 -> HClO}$ is $3 - 1 = 2e^-$.
$\Delta N$ for $\ce{HClO -> \tfrac{1}{2}Cl2}$ is $1 - 0 = 1e^-$.
Total electrons = $2 + 2 + 1 = 5e^-$ per Cl atom.
\begin{align*}
    5 \times E^\circ &= 2(1.181) + 2(1.674) + 1(1.630) \\
    &= 2.362 + 3.348 + 1.630 = 7.340\,\text{V} \\
    E^\circ &= \frac{7.340}{5} = 1.468\,\text{V} \approx +1.47\,\text{V}
\end{align*}
Hence option (a) is correct.
\end{solution}

\begin{problembox}
\textbf{CP.4.2} \hfill \textbf{[ATK / NRJ]}

From the Latimer diagram, does hypochlorous acid (\ce{HClO}) tend to undergo spontaneous disproportionation into \ce{HClO2} and \ce{Cl2} in $1\,\text{M}$ acid at $298\,\text{K}$?
\begin{enumerate}[label=(\alph*)]
    \item Yes, because $E^\circ_{\text{right}} (1.630\,\text{V}) < E^\circ_{\text{left}} (1.674\,\text{V})$.
    \item No, because $E^\circ_{\text{cell}} = 1.630 - 1.674 = -0.044\,\text{V} < 0$.
    \item Yes, because all chlorine oxyacids are universally explosive.
    \item No, because oxidation states cannot change without a catalyst.
\end{enumerate}
\end{problembox}
\begin{solution}
\textbf{(b)}
Disproportionation of \ce{HClO}:
\begin{align*}
    \text{Reduction: } & \ce{2HClO + 2H+ + 2e^- -> Cl2 + 2H2O} \quad (E^\circ = +1.630\,\text{V}) \\
    \text{Oxidation: } & \ce{HClO + H2O -> HClO2 + 2H+ + 2e^-} \quad (E^\circ = -1.674\,\text{V}) \\
    \text{Net: } & \ce{3HClO -> Cl2 + HClO2 + H2O}
\end{align*}
$E^\circ_{\text{disp}} = E^\circ_{\text{red}} - E^\circ_{\text{ox}} = 1.630 - 1.674 = -0.044\,\text{V}$.
Since $E^\circ_{\text{disp}} < 0$, $\Delta G^\circ > 0$, the reaction is \textbf{not} spontaneous under standard conditions. \ce{HClO} does not disproportionate into \ce{HClO2} and \ce{Cl2}.
\end{solution}

\begin{problembox}
\textbf{CP.4.3} \hfill \textbf{[ATK / NA]}

On a Frost diagram, the slope of the line segment joining two oxidation states represents:
\begin{enumerate}[label=(\alph*)]
    \item The equilibrium constant $K$ of the reaction.
    \item The standard reduction potential $E^\circ$ of the couple connecting those two states.
    \item The activation energy $E_a$ of the redox interconversion.
    \item The rate of electron transfer across the electrical double layer.
\end{enumerate}
\end{problembox}
\begin{solution}
\textbf{(b)}
By definition of the Frost diagram, the vertical axis is $y = N E^\circ = -\Delta G^\circ / F$ (relative to $N=0$).
The slope between two states $N_a$ and $N_b$ ($N_b > N_a$) is:
\begin{equation*}
    \text{Slope} = \frac{N_b E_b^\circ - N_a E_a^\circ}{N_b - N_a} = E^\circ(N_b / N_a)
\end{equation*}
Thus, the slope directly corresponds to the standard reduction potential of the redox couple. Steeper positive slope means stronger oxidizing ability.
\end{solution}
